\begin{textsample}{NEG Dim 2 – 2011\_c6 – Score -5.00 – 2011\_c6/enade\_2011\_tec\_autom\_c60e876b3b.txt}  \label{ex:f2_neg_025}
CAPÍTULO 3 ANÁLISE TÉCNICA DA PROVA Este capítulo tem por objetivo apresentar o desempenho dos estudantes concluintes de Tecnologia em Automação Industrial no ENADE/2011.
Para isso, foram calculadas as estatísticas básicas da prova em seu todo, bem como as estatísticas dos componentes relacionadas à Formação Geral, ao de Conhecimento Específico da Área e das questões discursivas isoladamente.
Nas tabelas, são apresentados o tamanho da população inscrita e de presentes, e as seguintes estatísticas das notas⁸ : média do desempenho na prova, erro padrão da média, desvio padrão, nota mínima, mediana e nota máxima.
As estatísticas apresentadas neste capítulo contemplam o total de estudantes concluintes da área de Tecnologia em Automação Industrial em 2011 do Brasil e, separadamente, por Grande Região.
Foram calculadas tendo -se em vista as seguintes agregações : ( a ) as Grandes Regiões e o país como um todo ; ( b ) a Categoria Administrativa ; e ( c ) a Organização Acadêmica.
Em relação aos gráficos de distribuição de notas, o intervalo considerado foi de 10 unidades, aberto à esquerda e fechado à direita, com exceção do primeiro intervalo, [ 0 ; 10 ], fechado em ambos os extremos.
Para os gráficos de distribuição das notas das questões discursivas, foram consideradas mais duas categorias : questão em branco e nota zero. 3.1 ESTATÍSTICAS BÁSICAS DA PROVA 3.1.1 Estatísticas Básicas Gerais A Tabela 3.1 apresenta as estatísticas básicas da prova por Grande Região.
A população total de inscritos foi de 1.959.
Destes, 1.577 estiveram presentes, sendo 19,5 \% o índice de não comparecimento.
A região de maior abstenção foi a Sudeste ( 22,1 \% ) e a de menor abstenção foi a Nordeste ( 11,9 \% ).
A média das notas da prova como um todo ( nas seções seguintes serão analisados os Componentes de Formação Geral e de Conhecimento Específico ) foi 40,5, sendo que os alunos da região Norte obtiveram a média mais baixa ( 35,1 ) e os da região Nordeste obtiveram a média mais alta ( 41,5 ).
As demais médias foram : 40,8 na região Sudeste, 40,3 na região Sul e 39,1 na região Centro-Oeste.
O desvio padrão para o Brasil como um todo foi 12,2, sendo o maior desvio padrão encontrado na região Sul ( 12,7 ) e o menor na região Nordeste ( 9,1 ), indicando uma menor dispersão das notas desta última região.
A região que obteve a maior nota máxima foi a Sudeste ( 80,6 ), ao passo que a região que atingiu a menor nota máxima foi a Norte ( 53,9 ).
A mediana do Brasil como um todo foi 40,5, sendo a maior mediana obtida na região Nordeste ( 41,4 ) e a menor obtida na Norte ( 32,5 ).
A nota mínima foi zero nas regiões Sudeste, Sul e Centro-Oeste.
Na região Norte a nota mínima foi 9,4, e na Nordeste, 21,2.
O comportamento das notas dos estudantes de todo o Brasil pode ser observado no Gráfico 3.1 que apresenta um histograma com a distribuição das mesmas.
Essa é uma distribuição unimodal com moda no intervalo ( 40 ; 50 ].
Apesar disso, nota -se que no intervalo ( 30 ; 40 ] a frequência de notas é bem próxima da encontrada no intervalo modal.
O coeficiente de assimetria da distribuição das notas é negativo ( – 0,11 ) e próximo de zero.
Por isso, pode -se considerar que a distribuição é aproximadamente simétrica.
Os coeficientes de assimetria das distribuições das regiões Centro-Oeste ( – 0,79 ), Norte e Sul ( – 0,24 para ambas ) são negativos, indicando concentração pouco maior do lado direito do histograma.
Já nas regiões Nordeste ( 0,20 ) e Sudeste ( 0,08 ), os coeficientes positivos e pequenos estão relacionados a uma discreta concentração à esquerda e cauda um pouco mais longa à direita. ⁸ Essas e outras estatísticas estão definidas no Capítulo 1. foi 12,2, sendo o maior desvio padrão encontrado na região Sul ( 12,7 ) e o menor na região Nordeste ( 9,1 ), indicando uma menor dispersão das notas desta última região.
A região que obteve a maior nota máxima foi a Sudeste ( 80,6 ), ao passo que a região que atingiu a menor nota máxima foi a Norte ( 53,9 ).
A mediana do Brasil como um todo foi 40,5, sendo a maior mediana obtida na região Nordeste ( 41,4 ) e a menor obtida na Norte ( 32,5 ).
A nota mínima foi zero nas regiões Sudeste, Sul e Centro-Oeste.
Na região Norte a nota mínima foi 9,4, e na Nordeste, 21,2.
O comportamento das notas dos estudantes de todo o Brasil pode ser observado no Gráfico 3.1 que apresenta um histograma com a distribuição das mesmas.
Essa é uma distribuição unimodal com moda no intervalo ( 40 ; 50 ].
Apesar disso, nota -se que no intervalo ( 30 ; 40 ] a frequência de notas é bem próxima da encontrada no intervalo modal.
O coeficiente de assimetria da distribuição das notas é negativo ( – 0,11 ) e próximo de zero.
Por isso, pode -se considerar que a distribuição é aproximadamente simétrica.
Os coeficientes de assimetria das distribuições das regiões Centro-Oeste ( – 0,79 ), Norte e Sul ( – 0,24 para ambas ) são negativos, indicando concentração pouco maior do lado direito do histograma.
Já nas regiões Nordeste ( 0,20 ) e Sudeste ( 0,08 ), os coeficientes positivos e pequenos estão relacionados a uma discreta concentração à esquerda e cauda um pouco mais longa à direita.
Os Gráficos 3.2, 3.3 e 3.4 apresentam informações referentes à média da nota final dos participantes, desagregando os resultados de acordo com, respectivamente, as Grandes Regiões do país, a Categoria Administrativa e a Organização Acadêmica.
Os gráficos apresentam o valor da média das notas como uma barra e os extremos do intervalo de confiança de 95 \% como linhas verticais unidas por uma linha horizontal na forma da letra H maiúscula.
Considerando -se o gráfico de notas segundo Grande Região ( Gráfico 3.2 ), bserva -se que não existe diferença estatisticamente significativa ao nível de 95 \% entre a maior média, obtida na região Nordeste ( 41,5 ) e a menor, obtida na região Norte ( 35,1 ).
Como apenas 16 alunos da região Norte realizaram a prova do ENADE/2011 ( ver tabela 2.3 ), o intervalo de confiança da média desta região é bem mais largo do que o das demais regiões.
Levando -se em conta os estudantes por Categorias Administrativas ( Gráfico 3.3 ), observa -se que existe diferença estatisticamente significativa entre as médias das notas das IES Públicas e Privadas.
Pode ser observado, ainda, que a média dos alunos de IES Públicas ( 43,3 ) é maior que a dos alunos de IES Privadas ( 39,1 ).
Tendo como foco o Gráfico 3.4, que apresenta as notas médias das provas segundo Organização Acadêmica, constata -se que a diferença entre as médias dos estudantes das Universidades ( 41,3 ) com relação aos alunos de Centros Universitários ( 39,0 ) é estatisticamente significativa ao nível de 95 \%. 3.1.2 Estatísticas Básicas no Componente de Formação Geral A Tabela 3.2 apresenta as estatísticas básicas em relação ao componente da prova que avalia a Formação Geral dos estudantes concluintes.
Os alunos de todo Brasil obtiveram desempenho médio de 47,8.
A maior média foi obtida na região Nordeste ( 53,1 ), e a menor, na região Centro-Oeste ( 45,3 ).
As demais médias foram : 46,8 na região Norte ; 48,9 na região Sudeste ; e 45,8 na região Sul.
Quanto à variabilidade, o desvio padrão das notas dos estudantes do Brasil como um todo foi 17,6.
Já o maior desvio padrão foi obtido na região Centro-Oeste ( 18,3 ) e o menor na região Norte ( 17,0 ).
Os demais desvios padrões foram : 17,8 na região Nordeste, 17,5 na região Sudeste e 17,3 na região Sul.
A maior nota no Componente de Formação Geral da prova do ENADE foi obtida por pelo menos um aluno da região Nordeste ( 100,0 ) enquanto que a menor nota máxima foi obtida na região Norte ( 82,5 ).
Nas outras regiões as notas máximas foram : 92,0 na região Sudeste, 86,5 na região Sul e 84,0 na região Centro-Oeste.
A mediana do Brasil como um todo foi 49,5, sendo a menor mediana encontrada na região Norte ( 43,3 ) e a maior encontrada na região Nordeste ( 53,8 ).
A nota mínima nesta parte foi zero nas regiões Sudeste, Sul e Centro-Oeste.
Na região Norte a nota mínima foi 25,0, e na região Nordeste, 15,0.
O Gráfico 3.5 propicia a avaliação do desempenho dos estudantes no Componente de Formação Geral a partir do histograma da distribuição das notas correspondentes.
A distribuição é unimodal, com moda em ( 50 ; 60 ], enquanto na prova como um todo a moda foi alcançada no intervalo ( 40 ; 50 ] com uma concentração no intervalo ( 30 ; 40 ] um pouco mais baixo.
Verifica -se, ainda, que no Gráfico 3.5 as notas apresentam maior dispersão do que no Gráfico 3.1 ( distribuição das notas da prova ).
Isso pode ser confirmado pela comparação dos desvios padrões : 12,2 para a nota da prova como um todo e 17,6 para o Componente de Formação Geral.
Para o Componente de Formação Geral, o coeficiente de assimetria da distribuição das notas dos estudantes, como na prova em sua totalidade, também é negativo ( – 0,29 ).
Por ser um coeficiente pequeno, a distribuição apresenta uma leve concentração à direita e cauda maior à esquerda.
Nas regiões Sudeste, Sul e Centro-Oeste os histogramas também possuem assimetria negativa ( entre – 0,29 e – 0,62 ).
Nas regiões Norte ( 0,73 ) e Nordeste ( 0,03 ) as distribuições possuem assimetria positiva : maior concentração de notas à esquerda e cauda mais longa conforme as notas aumentam, ao lado direito.
Nos Gráficos 3.6, 3.7 e 3.8 são apresentadas as informações referentes ao desempenho dos Concluintes no Componente de Formação Geral, em diferentes agregações : Grande Região do país, Categoria Administrativa e Organização Acadêmica.
Observa -se, pelo Gráfico 3.6, que existe diferença estatisticamente significativa entre as médias das notas no Componente de Formação Geral da região Nordeste, maior média, e das regiões Centro-Oeste e Sul.
O mesmo não ocorre em relação às regiões Sudeste e Norte.
Verifica -se, ainda, que o intervalo de confiança mais largo é o da região Norte, fato que está relacionado com o tamanho da população envolvida, menor na região Norte do que nas demais regiões.
No Gráfico 3.7, que representa as notas médias no Componente de Formação Geral segundo Categoria Administrativa do país, observa -se que existe diferença estatisticamente significativa entre as médias.
Os concluintes das IES Públicas obtiveram uma média maior ( 50,8 ) do que os das IES Privadas ( 46,4 ).
Considerando -se o tipo de Organização Acadêmica, nota -se, no Gráfico 3.8, que não há diferença estatisticamente significativa entre as médias das Universidades ( 49,1 ), Centros Universitários ( 47,5 ) e Faculdades ( 46,4 ). 3.1.3 Estatísticas Básicas do Componente de Conhecimento Específico A Tabela 3.3 apresenta as estatísticas básicas referentes ao Componente de Conhecimento Específico da área de Tecnologia em Automação Industrial.
A média do desempenho dos alunos do Brasil como um todo foi 38,0.
A maior média foi obtida na região Sul ( 38,5 ), e a menor, na região Norte ( 31,2 ).
As demais médias foram : 37,6 na região Nordeste, 38,1 na região Sudeste e 37,0 na região Centro-Oeste.
Quanto à variabilidade das notas, o desvio padrão do Brasil como um todo foi 13,4, sendo o maior desvio padrão observado na região Sul ( 14,2 ) e o menor na região Nordeste ( 10,1 ).
Os demais desvios foram : 13,6 na região Norte, 13,3 na região Sudeste e 13,0 na região Centro-Oeste.
Assim como os Gráficos 3.1 e 3.5, o Gráfico 3.9, apresentado a seguir, proporciona uma avaliação do desempenho de concluintes em relação ao Componente de Conhecimento Específico com um histograma da distribuição das notas correspondentes.
Dentre as três distribuições apresentadas, esta é a que apresenta maior concentração em notas baixas.
Esta também é uma distribuição unimodal, e o grupo modal é o ( 30 ; 40 ].
O coeficiente de assimetria da distribuição das notas do Componente de Conhecimento Específico é positivo e pequeno ( 0,13 ).
Nota -se, pelo histograma ( Gráfico 3.9 ), que esta é uma distribuição aproximadamente simétrica.
Já nas regiões Norte e Centro-Oeste os coeficientes de assimetria são negativos, iguais a – 0,29, evidenciando que a distribuição das notas dos estudantes dessas regiões, no Componente de Conhecimento Específico, tem cauda mais pesada à esquerda.
Na região Sul, apesar de negativo, o coeficiente é quase zero ( – 0,01 ).
Nas regiões Nordeste e Sudeste, este coeficiente é positivo ( 0,27 em ambas ) e maior do que o nacional.
A mediana das notas dos estudantes de todo o Brasil foi 37,4.
A maior mediana ocorreu na região Sul ( 39,4 ), e a menor, na região Norte ( 30,5 ).
As demais medianas foram : 36,7 na região Nordeste, 37,0 na região Sudeste e 35,4 na região Centro-Oeste.
A nota máxima do Brasil como um todo foi 85,9, sendo obtida por pelo menos um aluno da região Sudeste.
As demais notas máximas foram : 56,0 na região Norte, 64,4 na região Nordeste, 77,5 na região Sul e 69,8 na região Centro-Oeste.
A nota mínima foi zero em todas as regiões, com exceção da região Nordeste, cuja nota mínima foi 15,4.
Assim como os Gráficos 3.1 e 3.5, o Gráfico 3.9, apresentado a seguir, proporciona uma avaliação do desempenho de concluintes em relação ao Componente de Conhecimento Específico com um histograma da distribuição das notas correspondentes.
Dentre as três distribuições apresentadas, esta é a que apresenta maior concentração em notas baixas.
Esta também é uma distribuição unimodal, e o grupo modal é o ( 30 ; 40 ].
O coeficiente de assimetria da distribuição das notas do Componente de Conhecimento Específico é positivo e pequeno ( 0,13 ).
Nota -se, pelo histograma ( Gráfico 3.9 ), que esta é uma distribuição aproximadamente simétrica.
Já nas regiões Norte e Centro-Oeste os coeficientes de assimetria são negativos, iguais a – 0,29, evidenciando que a distribuição das notas dos estudantes dessas regiões, no Componente de Conhecimento Específico, tem cauda mais pesada à esquerda.
Na região Sul, apesar de negativo, o coeficiente é quase zero ( – 0,01 ).
Nas regiões Nordeste e Sudeste, este coeficiente é positivo ( 0,27 em ambas ) e maior do que o nacional.
Os Gráficos 3.10, 3.11 e 3.12 apresentam uma comparação dos resultados em relação à Grande Região do país, à Categoria Administrativa e à Organização Acadêmica, agora levando em conta o desempenho de alunos no Componente de Conhecimento Específico da prova.
Pelo Gráfico 3.10, observa -se que não há diferença estatisticamente significativa entre as médias das notas no Componente de Conhecimento Específico das regiões.
Quanto à Categoria Administrativa ( gráfico 3.11 ), constata -se um comportamento semelhante àquele da parte de Formação Geral e à prova como um todo, ou seja, existe diferença estatisticamente significativa entre as médias das IES Públicas ( 40,8 ) e Privadas ( 36,7 ), sendo que a maior média foi obtida por alunos de IES Públicas de ensino.
Quanto ao Gráfico 3.12, observa -se que entre as Faculdades ( 38,3 ) e Universidades ( 38,7 ) não existe diferença estatisticamente significativa das notas no Componente de Conhecimento Específico.
Já a diferença das notas entre os Centros Universitários ( 36,2 ) e Universidades é significativa, sendo a média dos Concluintes das Universidades maior do que a dos Centros Universitários. 3.2 ANÁLISE DAS QUESTÕES OBJETIVAS 3.2.1 Componente de Formação Geral A Tabela 3.4 apresenta as estatísticas básicas relativas às oito questões objetivas do componente da prova que abrange a Formação Geral dos estudantes.
A média do Brasil foi 48,8.
A menor média foi encontrada na região Centro-Oeste ( 45,8 ), e a maior, na região Nordeste ( 53,9 ).
As demais médias foram : 48,4 na região Norte, 49,4 na região Sudeste e 47,6 na região Sul.
O desvio padrão do Brasil foi 19,4, sendo o maior desvio padrão encontrado na região Nordeste ( 22,3 ) e o menor na região Sul ( 18,8 ).
Os demais desvios foram : 21,3 na região Nordeste, 19,4 na região Sudeste e 20,0 na região Centro-Oeste.
A mediana das notas de todo o Brasil ( 50,0 ) só não foi a mesma na região Norte, onde a mediana foi 43,8.
A nota máxima foi 100,0 nas regiões Nordeste e Sudeste, mas nas demais a nota máxima atingida foi 87,5.
A nota mínima foi diferente de zero apenas na região Norte ( 25,0 ).
A Tabela 3.5 apresenta os índices de facilidade e discriminação ( ponto bisserial ) das questões objetivas do Componente de Formação Geral.
Quanto ao índice de facilidade, foram usadas as seguintes cores para diferenciar o nível de dificuldade da questão : • Azul para as questões classificadas com índice muito fácil ( índice \&gt ; =0,86 ), verde para as questões classificadas com índice fácil ( 0,61 a 0,85 ), amarelo para as questões classificadas com médio ( 0,41 a 0,60 ), vermelho para as questões classificadas com difícil ( 0,16 a 0,40 ) e roxo para as questões classificadas com muito difícil ( \&lt ; =0,15 ).
Já quanto ao índice de discriminação, foram usadas as seguintes cores para qualificar a questão : • As questões classificadas com índice fraco receberam a cor vermelho ( índice \&lt ; =0,19 ), as classificadas com médio receberam a cor amarelo ( 0,20 a 0,29 ), as classificadas com bom receberam a cor verde ( 0,30 a 0,39 ) e as classificadas com muito bom ( \&gt ; =0,40 ) receberam a cor azul.
As questões objetivas do Componente de Formação Geral, segundo o índice de facilidade, foram assim avaliadas : das oito questões, nenhuma teve o índice de facilidade classificado como muito fácil.
Quatro questões foram tidas como fáceis, por terem índice de acertos situado na faixa entre 0,61 e 0,85 ( de 61,0 \% a 85,0 \% de acertos ).
Uma questão foi considerada de dificuldade média, situando -se no intervalo entre 0,41 e 0,60 do índice de facilidade, ou seja, houve entre 41,0 \% e 60,0 \% de acertos.
Três questões foram classificadas como difíceis, situando -se no intervalo entre 0,16 e 0,40.
Por fim, nenhuma questão foi classificada como muito difícil por apresentar menos de 15 \% de acertos.
A Tabela 3.5 apresenta os índices de facilidade e discriminação ( ponto bisserial ) das questões objetivas do Componente de Formação Geral.
Quanto ao índice de facilidade, foram usadas as seguintes cores para diferenciar o nível de dificuldade da questão :  Azul para as questões classificadas com índice muito fácil ( índice \&gt ; =0,86 ), verde para as questões classificadas com índice fácil ( 0,61 a 0,85 ), amarelo para as questões classificadas com médio ( 0,41 a 0,60 ), vermelho para as questões classificadas com difícil ( 0,16 a 0,40 ) e roxo para as questões classificadas com muito difícil ( \&lt ; =0,15 ).
Já quanto ao índice de discriminação, foram usadas as seguintes cores para qualificar a questão :  As questões classificadas com índice fraco receberam a cor vermelho ( índice \&lt ; =0,19 ), as classificadas com médio receberam a cor amarelo ( 0,20 a 0,29 ), as classificadas com bom receberam a cor verde ( 0,30 a 0,39 ) e as classificadas com muito bom ( \&gt ; =0,40 ) receberam a cor azul.
As questões objetivas do Componente de Formação Geral, segundo o índice de facilidade, foram assim avaliadas : das oito questões, nenhuma teve o índice de facilidade classificado como muito fácil.
Quatro questões foram tidas como fáceis, por terem índice de acertos situado na faixa entre 0,61 e 0,85 ( de 61,0 \% a 85,0 \% de acertos ).
Uma questão foi considerada de dificuldade média, situando -se no intervalo entre 0,41 e 0,60 do índice de facilidade, ou seja, houve entre 41,0 \% e 60,0 \% de acertos.
Três questões foram classificadas como difíceis, situando -se no intervalo entre 0,16 e 0,40.
Por fim, nenhuma questão foi classificada como muito difícil por apresentar menos de 15 \% de acertos.
O Gráfico 3.13, para exemplificar, analisa o comportamento da questão de número 3 de Formação Geral.
Trata -se da questão classificada como fácil e que obteve o maior índice de discriminação dessa parte da prova.
Como já comentado, para análise das questões objetivas relativas à Formação Geral segundo o poder de discriminação, utilizou -se, o índice de discriminação ponto bisserial.
Nesta análise as questões foram assim avaliadas : cinco das oito questões apresentaram índices acima de 0,40 e, assim, foram classificadas com índice muito bom para esse grupo de alunos ; duas questões tiveram bom índice de discriminação, entre 0,30 e 0,39.
Apenas uma foi classificada com índice médio e nenhuma teve índice fraco.
O índice de facilidade variou de 0,16 a 0,75, e o de discriminação, de 0,26 a 0,53.
As questões com índices de discriminação muito bom, de números 1, 2, 3, 4 e 5, se classificam quanto ao índice de facilidade da seguinte forma : três na categoria fácil ( questões 1, 3 e 5 ), uma na categoria médio ( questões 4 ) e uma difícil ( questão 2 ).
Em particular, a questão 3 foi a que apresentou maior poder discriminatório, com índice 0,53, e foi uma das classificadas como fácil, com uma proporção de 0,64 de acertos.
A questão de número 8 apresentou menor índice de facilidade 0,16, ou seja, um quantitativo pequeno de estudantes conseguiu resolvê -la, dentro do universo de participantes.
O índice de discriminação desta questão foi médio.
Já a questão de número 7, também considerada difícil, obteve índice de discriminação 0,33, considerado bom.
O Gráfico 3.13, para exemplificar, analisa o comportamento da questão de número 3 de Formação Geral.
Trata -se da questão classificada como fácil e que obteve o maior índice de discriminação dessa parte da prova.
Neste gráfico, cada uma das cinco curvas representa o percentual de respostas em determinada alternativa da questão, em função do número de acertos dos estudantes nesta parte da prova ( Formação Geral/Múltipla Escolha ), antes de possíveis eliminações pelo critério do ponto bisserial.
A curva em vermelho corresponde à alternativa E, a correta para esta questão.
Assim, observa -se que entre os estudantes com notas mais baixas, nessa parte do exame, a situação mais frequente foi a escolha de uma das alternativas incorretas : a alternativa A ( em azul ).
Na medida em que o número de acertos aumenta, indicando desempenho melhor nesta parte da prova, aumenta concomitantemente a proporção de estudantes que selecionaram a alternativa correta E, atingindo 100 \% para os que acertaram as oito questões.
Essa análise permite verificar como a questão discriminou os grupos de desempenho, justificando o alto índice obtido na questão.
Os gráficos relativos às demais questões de Formação Geral constam do Anexo I. 3.2.2 Componente de Conhecimento Específico A Tabela 3.6 apresenta as estatísticas básicas em relação às questões objetivas do Componente de Conhecimento Específico da prova, por Grande Região.
A média do Brasil deste componente foi de 40,3.
A menor média foi observada na região Norte ( 32,7 ), e a maior, na região Sul ( 40,8 ).
O desvio padrão de todo o Brasil foi 14,1, sendo o menor desvio padrão encontrado na região Nordeste ( 11,2 ) e o maior na região Sul ( 14,9 ).
A mediana de todo o Brasil foi 38,1, a mesma das regiões Nordeste, Sudeste e Centro-Oeste.
Na região Sul verifica -se a maior mediana ( 42,9 ), e a menor na região Norte ( 31,0 ).
A nota máxima das questões objetivas do Componente de Conhecimento Específico foi 90,5, obtida por pelo menos um aluno da região Sudeste.
Nas demais regiões a nota máxima alcançada foi : 57,1 na região Norte, 71,4 na região Nordeste, 76,2 nas regiões Sul e Centro-Oeste.
A nota mínima foi 0,0 em todas as regiões, exceto na região Nordeste ( 14,3 ).
A Tabela 3.7 apresenta os índices de facilidade e discriminação ( ponto bisserial ) das questões objetivas do Componente de Conhecimento Específico.
Para facilitar a diferenciação das questões usou -se as mesmas cores da Tabela 3.5 para as diferentes classificações dos índices de facilidade e de discriminação.
Dentre as questões objetivas da parte da prova relativa ao Componente de Conhecimento Específico, nenhuma foi anulada pela Comissão.
Desse modo, a classificação quanto ao índice de facilidade foi estabelecida com base nas 27 questões.
A partir dos índices obtidos, pode -se concluir que a maioria das questões objetivas da prova foi considerada pelo menos difícil : das 27 questões, quatro foram classificadas como muito difíceis e 14 como difíceis.
Não houve questão classificada como muito fácil, ao passo que quatro foram tidas como fáceis, na faixa de 0,61 a 0,85 do índice de facilidade, e outras cinco consideradas médias, entre 0,41 e 0,60.
Já quanto aos índices de discriminação das questões objetivas do Componente de Conhecimento Específico da prova, tem -se como resultado a seguinte classificação : oito das 27 questões foram consideradas como boas, enquanto apenas duas tiveram índice de discriminação muito bom.
Assim, para dez das 27 questões os índices de discriminação foram bons ou muito bons.
Dentre as demais, onze delas foram classificadas como médias e outras seis como fracas, sendo dezessete, por conseguinte, a quantidade de questões nos dois patamares mais baixos de discriminação.
Constata -se, assim, que a prova – no que se refere ao Componente de Conhecimento Específico – possuía capacidade baixa de discriminar entre aqueles que dominam ou não o conteúdo.
Apenas duas questões alcançaram os maiores índices de discriminação, as de números 20 e 33 classificadas com índice muito bom, situando -se no intervalo de 0,40 a 0,45 do índice.
Quanto ao índice de facilidade uma delas ( questão 20 ) foi classificada na categoria média e a outra ( questão 33 ) na categoria fácil.
A questão de número 25 foi a mais difícil dentre as 27 questões específicas, com baixo índice de facilidade, apenas 6,0 \% de acertos.
Essa questão apresentou poder discriminatório igualmente baixo, 0,17 ( fraco ), o que comprova ter sido esta a mais difícil para os estudantes.
Destacam -se, também, as questões 9 e 31 classificadas como muito difíceis.
A questão 9, com 7,0 \% de percentual de acerto teve índice de facilidade fraco ( 0,10 ) e o mesmo ocorreu com a questão 31, muito difícil ( 10,0 \% de acerto ) e fraco índice de discriminação ( 0,09 ).
Por isso, as questões 9, 25 e 31 foram eliminadas do cômputo da nota final.
Além destas três, as demais questões com índice fraco de discriminação, questões 24, 26 e 27 também não foram computadas.
A título de exemplo das análises do comportamento das questões objetivas, o Gráfico 3.14 analisa a questão 20 do Componente de Conhecimento Específico.
Esta foi a questão que obteve maior índice de discriminação ( 0,42 ) e foi classificada com índice de facilidade médio 0,57, ou seja, 57,0 \% dos estudantes assinalaram acertadamente a opção A, correspondente ao gabarito.
Neste gráfico, cada uma das cinco curvas representa o percentual de respostas em determinada alternativa da questão 20, em função da quantidade de acertos dos estudantes nesta parte da prova, antes de possíveis eliminações de questões pelo critério do ponto bisserial.
A alternativa correta A, representada no gráfico pela curva em azul, foi escolhida em maiores proporções pelos alunos com desempenho melhor nesta parte da prova.
Já as alternativas incorretas, também denominadas distratores, foram selecionadas principalmente por aqueles com notas mais baixas.
Aqueles com nota zero, na sua quase totalidade deixaram esta questão em branco ou marcaram mais de uma alternativa, comportamento considerado inválido.
A proporção de alunos que selecionou a resposta correta A aumenta gradativamente, chegando a atingir 100 \% para os alunos com 17 acertos ou mais, enquanto a proporção dos que escolheram alternativas incorretas, por exemplo as alternativas : B – linha verde, C – linha preta e E – linha vermelha, decai a partir de um acerto, como função do número de acertos nesta parte da prova.
Os gráficos relativos às demais questões do Componente de Conhecimento Específico constam do Anexo I. 3.3 ANÁLISE DAS QUESTÕES DISCURSIVAS 3.3.1 Componente de Formação Geral As análises dos resultados de desempenho dos estudantes de Tecnologia em Automação Industrial nas duas questões discursivas relativas à Formação Geral encontram -se na Tabela 3.8 e no Gráfico 3.15.
Na tabela 3.8 observa -se que os estudantes de todo o Brasil obtiveram média 46,3, nas questões discursivas de Formação Geral, pouco menor do que nas questões objetivas, para as quais a média foi 48,8.
A mediana confirma um desempenho semelhante dos alunos de todo o Brasil nos dois tipos de questões do Componente de Formação Geral.
Como consta da tabela 3.8, esta estatística foi de 48,8 para questões discursivas, e nas questões objetivas ( tabela 3.4 ) essa estatística foi ligeiramente maior, 50,0.
Já em relação ao desvio padrão, pode -se notar um aumento : 19,4, nas questões objetivas do Componente de Formação Geral dos alunos de todo o Brasil, para 26,8 nas questões discursivas do mesmo componente.
A maior mediana foi encontrada na região Nordeste ( 55,0 ), e a menor, na região Norte ( 42,5 ).
A nota máxima ( 100,0 ) foi obtida por alunos das regiões Nordeste, Sudeste e Sul.
Na região Norte a nota máxima foi 85,0, e na região Centro-Oeste foi 97,5.
A nota mínima ( 0,0 ) foi a mesma em todas as regiões.
O Gráfico 3.15 representa a distribuição das notas nas questões discursivas no Componente de Formação Geral.
A moda desta distribuição ocorre no intervalo ( 40 ; 50 ].
Observa -se outro intervalo de frequência destacada é o [ 0 ; 10 ], provavelmente devido às notas zero e aos alunos que deixaram este tipo de questão em branco.
O coeficiente de assimetria negativo desta distribuição ( – 0,24 ) indica assimetria à esquerda, com maior concentração em notas mais baixas.
A assimetria é negativa, e da mesma ordem de \textbf{grandeza} do que a obtida para o Brasil, em todas as regiões variando de -0,29 a -0,14.
A seguir, serão analisados os desempenhos da Área de Tecnologia em Automação Industrial nas duas questões discursivas de Formação Geral do ENADE/2011, comparando os resultados obtidos com comentários para cada questão.
Os comentários da Banca de docentes corretores a respeito do observado na correção das respostas dos estudantes, suas impressões e conclusões serão apresentados junto à análise de cada questão.
Cumpre esclarecer que, tendo em vista que as questões discursivas de Formação Geral são padronizadas, ou seja, constam de todas as provas, os comentários da Banca são os mesmos para todas as carreiras acadêmicas, sendo direcionados a todos os estudantes que participaram do ENADE/2011. 3.3.1.1 Análise da Questão Discursiva 1 do Componente de Formação Geral Os dados de Tecnologia em Automação Industrial, obtidos a partir das respostas à questão 1, encontram -se na Tabela 3.9 e no Gráfico 3.16.
Nessa questão, de melhor desempenho dentre as duas discursivas de Formação Geral, os alunos de todo Brasil tiveram média 54,2.
A maior média para a questão 1 foi obtida na região Nordeste ( 56,3 ), e a menor, na região Norte ( 48,1 ).
Quanto à variabilidade das notas, o desvio padrão de todo o Brasil foi 34,3.
O menor desvio padrão foi obtido na região Sudeste ( 33,7 ), e o maior, na região Norte ( 38,9 ).
As notas máximas e mínimas da questão discursiva 1 foram as mesmas para todas as regiões do Brasil, respectivamente, 0,0 e 100,0.
A mediana de todo o país, como também das regiões Nordeste e Sudeste, foi 60,0.
A menor mediana foi observada na região Norte ( 30,0 ). e a maior, na região Centro-Oeste ( 65,0 ) O Gráfico 3.16 mostra a distribuição das notas na questão discursiva 1 do Componente de Formação Geral.
Observa -se que a moda ocorre no intervalo ( 90 ; 100 ], indicando que em torno de 18 \% dos concluintes da Área de Tecnologia em Automação Industrial não tiveram dificuldade com esta questão.
O coeficiente de assimetria da distribuição das notas de todo o Brasil nesta questão é negativo ( – 0,15 ).
O mesmo ocorre nas demais regiões ( entre – 0,33 e – 0,06 ), exceto na região Norte ( 0,25 ). 3.3.1.2 Comentários sobre a correção das respostas à Questão Discursiva 1 De maneira geral, a aplicação da chave de correção da questão 1 de Formação Geral não apresentou qualquer dificuldade digna de menção.
Para isso contribuíram, indubitavelmente, os ajustes feitos após a correção da amostra e a reunião entre todos os membros da Banca de docentes corretores.
As poucas dúvidas, todas pontuais, apresentadas pelos corretores, foram acompanhadas e respondidas pela coordenação e subcoordenação da correção das questões de Formação Geral, por meio da ferramenta de Gerenciamento de Dúvidas do Sistema de Correção On-line.
Não houve registro de qualquer ocorrência que pusesse em xeque o padrão de resposta ou a efetividade e a adequação da chave de correção.
Explica -se : trata -se de questão com comando claro, direto e objetivo ( solicitava -se, basicamente, três vantagens justificadas de cursos a distância ), cujas respostas foram corrigidas por meio da aplicação de um chave de correção testada e aprovada previamente.
Havia absoluta clareza quanto aos critérios de avaliação da correspondência entre as respostas dos estudantes e as possibilidades de vantagens de cursos a distância admitidas como corretas no padrão de resposta oficial, além de gradações explícitas ( e fáceis de aplicar ) dos diferentes níveis de pontuação previstos.
Felizmente, portanto, não há reparo a registrar em relação à facilidade de aplicação do padrão de resposta e da chave de correção, e nem em relação à atribuição dos diferentes níveis de pontuação previstos.
Todas as dificuldades que poderiam ter obstado a correta aplicação do padrão de resposta oficial e da respectiva chave de correção foram evitadas por meio dos ajustes feitos após a correção da amostra e fartamente debatidos com toda a Banca.
Digno de nota é que quantidade tão significativa de profissionais envolvidos na mesma tarefa – tanto para a questão 1 quanto para a 2 – tenha apresentado tão poucas dificuldades na execução da correção, em termos proporcionais.
Em suma, a correção da questão 1 da prova de Formação Geral do ENADE 2011 foi exemplarmente bem planejada, servindo -se de padrão de resposta muito bem adequado à questão proposta.
Quanto ao tema desta questão, em particular, a Banca verificou que uma parcela significativa de estudantes evocou experiências bastante concretas e próximas de sua realidade.
Houve várias respostas que indicavam uma vivência pessoal de ensino superior na modalidade Educação a Distância ( EaD ), evidenciando o tom de depoentes nos textos apresentados.
Os estudantes, em sua grande maioria, utilizaram parte considerável do espaço de 15 linhas disponíveis para a resposta – e outra parcela menos significativa dos que não o fizeram demonstraram notável capacidade de atender ao comando da questão de maneira objetiva, curta e, via de regra, correta.
Registre -se, ainda, que foram relativamente poucos os casos de respostas que tenham passado ao largo do tema em pauta na questão.
As capacidades de leitura, de compreensão do comando proposto e de expressão escrita que os estudantes avaliados na edição 2011 do ENADE foram satisfatórias.
Não obstante, seria leviano perder de vista que a qualidade dos textos redigidos em resposta às questões discursivas do Exame ainda está muito aquém do que se espera de concluintes de cursos de ensino superior de todas as regiões do país.
Quanto ao conteúdo das respostas, a Banca constatou boa capacidade, por parte da maioria dos estudantes, de compreensão do tema e do comando da questão.
Foram relativamente poucos os casos de respostas que deixaram de enumerar vantagens da modalidade EaD, e proporcionalmente escassos os estudantes que citaram vantagens não previstas no padrão de resposta.
Foi frequente, a tentativa direta de atender ao comando da questão.
Os erros mais comuns, em relação ao padrão de respostas e à grade de correção, foram fruto do desdobramento em vários “ itens ” daquilo que, de acordo com o padrão de respostas oficial, representava uma única vantagem.
Destacaram -se, neste caso, as respostas que apontavam a flexibilidade de horário e/ou local como duas vantagens distintas daquela modalidade de ensino.
Quanto aos diferentes níveis de pontuação previstos, a maior causa de baixas pontuações foi a ausência de justificativas, e mesmo de argumentação, para uma ou mais das vantagens enumeradas.
Isso demonstra não apenas a objetividade das respostas, coerente com a objetividade do comando da questão ( “ enumere três vantagens de um curso a distância ” ), mas também certa dificuldade de formulação plena de um texto, ou ao menos de parágrafos, em formato dissertativo – mesmo diante de uma média de 5 linhas disponíveis para cada vantagem a enumerar.
Os acertos mais comuns, ou seja, os “ itens ” do padrão de respostas mais frequentemente mencionados foram : ( 1 ) a flexibilidade de horário e/ou local ; ( 2 ) a capilaridade do ensino a distância ; ( 3 ) a democratização do acesso à educação de qualidade ; e ( 4 ) os custos menores que os de cursos presenciais.
Dentre as vantagens previstas no padrão de respostas que foram menos citadas, destacam -se a inclusão de pessoas com comprometimento motor, a qualificação de professores e a troca de experiências entre os participantes.
Foram muito comuns, no entanto, as menções à supostamente maior facilidade de acesso a professores ou/e tutores em cursos superiores a distância. 3.3.1.3 Análise da Questão Discursiva 2 do Componente de Formação Geral A Tabela 3.10 mostra que o desempenho dos estudantes na questão 2 ( média 38,5 ) foi inferior ao obtido na questão de número 1 ( média 54,2 ).
A região Nordeste foi aquela onde a média, nessa questão, foi maior ( 47,7 ), e a de menor média foi a região Centro-Oeste ( 33,3 ).
Quanto à variabilidade das notas, o desvio padrão de todo o Brasil foi 30,4, Quanto ao desvio nessa questão em todas as regiões ficou em torno de 30,3, exceto na região Norte ( 23,3 ), o menor obtido.
A mediana das notas das regiões foi 40,0 ( Sul e Centro-Oeste ) ou 50,0 ( Norte, Nordeste e Sudeste ).
A nota máxima ( 100,0 ) foi obtida em todas as regiões, com exceção da região Norte ( 70,0 ).
A nota mínima foi 0,0 em todas as regiões, sem exceção.
O Gráfico 3.17 mostra a distribuição das notas na questão discursiva 2 do Componente de Formação Geral.
Este gráfico não apresenta apenas uma moda, mas possui dois máximos locais.
O intervalo modal, o de maior frequência, é aquele que representa as questões deixadas em branco, com mais de 25 \% dos estudantes.
Dentre as notas daqueles que resolveram a questão, a moda é o intervalo ( 50 ; 60 ].
O coeficiente de assimetria, apesar de negativo é bem próximo de zero ( – 0,02 ), indicando uma distribuição aproximadamente simétrica.
Os coeficientes de assimetria das regiões Sudeste ( – 0,10 ), Nordeste ( – 0,37 ) e Norte ( – 0,83 ) também são negativos, e nesta última região indica uma assimetria à esquerda um pouco mais acentuada.
Nas regiões Sul e Centro-Oeste as distribuições possuem assimetria positiva com coeficiente pequeno, 0,14 e 0,17, respectivamente. 3.3.1.4 Comentários sobre a correção das respostas à Questão Discursiva 2 Diferentemente da questão 1, a questão 2, cujo tema aborda políticas públicas para a erradicação do analfabetismo e para o aumento da empregabilidade, permitiu mais amplitude na elaboração das respostas dos estudantes, uma vez que solicita a apresentação de uma proposta.
À falta de objetividade técnica na elaboração da questão – sobretudo no seu comando – correspondeu uma miríade de possibilidades interpretativas.
Registre -se que não raras foram as respostas que se utilizaram de lugares comuns e exortações religiosas/humanitárias/cívicas, contudo, na grande maioria dos casos, essas exortações foram usadas como um encerramento do texto e não comprometeram a resposta.
Muito frequente também foi a confusão entre política pública e política partidária, bem como alguns poucos, confundiam programas educacionais com programas educativos nas TVs.
Assim, destacaram -se como propostas/programas mais recorrentes : • prosseguimento das políticas já existentes, mas com o aumento dos investimentos, normalmente sugerindo a ampliação da bolsa família, da Educação de Jovens e Adultos ( EJA ) e da EaD ; • valorização do magistério, construção de escolas e melhoria das atuais ; • parceria do governo com empresas para manutenção de salas de aula visando a alfabetização de seus funcionários, oferecimento de estágios e redução da carga horária em troca de \textbf{isenção} de impostos ; • parcerias com igrejas e ONGs para criação de espaços de alfabetização ; • escolas itinerantes e alfabetizadores em domicílio, principalmente para pessoas com dificuldade de locomoção, como os idosos e deficientes físicos ; • erradicação do trabalho infantil ; • vinculação da bolsa família não apenas à frequência, mas também e, principalmente, ao resultado obtido pelo aluno na escola ; • revisão das políticas atuais, sendo a mais frequente a extinção da aprovação automática ; • críticas consistentes em relação a modelos didáticos considerados inadequados e desestimulantes para a educação não só de adultos, mas de pessoas de todas as idades.
Destacam -se, ainda, outras sugestões apresentadas : • \textbf{Creches} nas escolas onde os pais estão sendo alfabetizados. • Diminuição de duas horas na jornada de trabalho em empresas para funcionários não alfabetizados, para que possam frequentar a escola. • Campanhas educativas vinculadas aos meios de comunicação. • Aumento do número de escolas noturnas. • Formação específica para professores alfabetizadores.
Quanto à relação entre o analfabetismo e a empregabilidade, deve -se sublinhar que nem todos os estudantes estabeleceram claramente o vínculo entre essas duas situações sociais.
Alguns falaram separadamente de uma e de outra.
Mas a maioria fez referência à necessidade de estudo para “ conseguir um bom emprego com um bom salário ”.
Alguns, em menor número, estabeleceram de forma bastante interessante a questão histórica para a situação do Nordeste ; e, também, a relação entre a pessoa analfabeta/com pouca instrução e a desempregada/empregada em serviços mais pesados e pior remunerados, que não teriam condições de educar seus filhos que, por sua vez, também não teriam melhores oportunidades no mercado de trabalho, identificando a formação de um círculo \textbf{vicioso} e a necessidade de sua interrupção.
Muitos afirmaram que o analfabetismo não é o único responsável pelo desemprego e sim a má distribuição de renda.
Grande também foi o número de alunos que criticaram o resultado da pesquisa, afirmando que a mesma é enganosa, uma vez que considera alfabetizada a pessoa que “ desenha seu nome ”, sendo comum associarem a atual política de alfabetização com ganhos eleitorais.
Foi comum a resposta incompleta, em que o estudante fez apenas a análise das desigualdades/crítica do quadro apresentado, ou só apresentou proposta.
Alguns se limitaram a responsabilizar o governo referindo -se de forma bastante genérica à questão da “ educação ” e “ profissional ”.
Entre os equívocos que mais se repetiram, destacou -se a simples análise dos dados apresentados na tabela que consta do enunciado.
Alguns textos, inclusive, estavam corretos, sem que, no entanto, fosse respondida a pergunta.
Da mesma forma, quando os estudantes partiram para a segunda parte da questão, fizeram referência a vários programas já existentes ou simplesmente disseram que é muito importante que existam projetos para a educação.
Outro equívoco recorrente foi a análise da educação no Brasil como um todo.
Essa análise, apesar de correta, não se referia especificamente à questão do analfabetismo, o que, por conseguinte, levava a sugestões que não eram direcionadas à erradicação do analfabetismo, e sim à melhora da educação no Brasil.
Sendo assim, foram apresentadas propostas como o aumento do número de faculdades, o reforço de alunos do Ensino Médio, a criação de escolas técnicas, etc.
Por outro lado, houve um grande número de redações bem escritas e precisas, no que se refere ao que foi exigido pela questão.
Foram análises equilibradas e sensatas, correlacionando a problemática do analfabetismo com o desemprego e a apresentação de sugestões bastante consistentes. 3.3.2 Componente de Conhecimento Específico Na parte da prova relativa às questões discursivas no Componente de Conhecimento Específico ( Tabela 3.11 ), observa -se que a média ( 25,2 ) foi mais baixa do que para as questões objetivas deste componente ( 40,3 ).
Comparando -se estes dados com os das questões discursivas do Componente de Formação Geral, também se observa um resultado inferior.
Enquanto na parte de Conhecimento Específico a média foi 25,2, no Componente de Formação Geral a média das questões discursivas para estudantes de Tecnologia em Automação Industrial de todo o Brasil foi 46,3.
A maior média deste componente foi obtida pelos estudantes da região Sul ( 25,4 ), e a menor pelos da região Norte ( 22,7 ).
Quanto à variabilidade das notas, o desvio padrão de todo o Brasil foi 18,0.
O maior desvio padrão foi encontrado na região Sul ( 19,3 ), e o menor, na região Nordeste ( 15,7 ).
A maior nota máxima foi obtida na região Sul ( 85,0 ), enquanto a menor nota máxima foi encontrada na região Norte ( 51,7 ).
Além disso, a nota mínima ( 0,0 ) foi obtida por estudantes de todas as regiões do Brasil, sem exceção.
A mediana do país como um todo foi 25,0, também encontrada nas regiões Sudeste e Sul.
As demais medianas foram : 21,7 na região Norte, 23,3 na região Nordeste e 26,7 na região Centro-Oeste.
O Gráfico 3.18 representa a distribuição das notas nas questões discursivas no Componente de Conhecimento Específico.
Essa distribuição é unimodal com moda no intervalo de [ 0 ; 10 ].
Registra -se uma tendência decrescente a partir do intervalo ( 20 ; 30 ] até o intervalo ( 80 ; 90 ].
Não foram registradas notas maiores que 85,0.
O coeficiente de assimetria da distribuição das médias das notas nas três questões discursivas do Componente de Conhecimento Específico é 0,27, relacionado à longa cauda à direita.
Para todas as Grandes Regiões o coeficiente também é positivo, variando entre 0,09 e 0,33. 3.3.2.1 Análise da Questão Discursiva 3 do Componente de Conhecimento Específico Na questão 3, cujos resultados aferidos encontram -se descritos na Tabela 3.12, a média dos estudantes de todo o Brasil foi 38,5.
A menor média nessa questão foi obtida pelos alunos da região Norte ( 34,1 ), enquanto a maior média foi obtida na região Nordeste ( 41,4 ).
Quanto à variabilidade das notas, o desvio padrão de todo o Brasil foi 26,7.
O maior desvio padrão foi obtido na região Sul ( 27,8 ), enquanto o menor foi obtido na região Nordeste ( 23,3 ).
A nota máxima, 100,0 pontos, foi alcançada por pelo menos um aluno das regiões Sudeste e Sul.
A região onde ocorreu a menor nota máxima foi a Norte ( 70,0 ).
A mediana do Brasil como um todo foi 40,0, a mesma das regiões Sudeste e Sul.
Nas regiões Norte e Centro-Oeste a mediana foi 45,0, e na região Nordeste, 42,5.
A nota mínima ( 0,0 ) foi obtida em todas as regiões do Brasil, sem exceção.
O Gráfico 3.19 mostra a distribuição das notas na questão discursiva 3, do Componente de Conhecimento Específico.
Essa distribuição é bimodal com moda nas questões em branco e no intervalo ( 40 ; 50 ].
O coeficiente de assimetria, apesar de negativo, é bem pequeno ( – 0,04 ), indicando uma distribuição aproximadamente simétrica.
Nas Grandes Regiões as distribuições possuem coeficiente de assimetria também negativo e próximo de zero ( entre – 0,27 e – 0,04 ), a menos da região Sul cujo coeficiente é 0,04. 3.3.2.2 Comentários sobre a correção das respostas à Questão Discursiva 3 Questão aborda conhecimento geral na área técnico-científica, tratando de itens sobre segurança do trabalho, aplicados especificamente nos serviços em eletricidade.
Simples, bem elaborada, exigindo que o estudante conheça a terminologia e equipamentos típicos, as normas regulamentadoras e os procedimentos empregados.
A pontuação dos itens foi adequada à dificuldade.
A questão como um todo foi considerada fácil, porém não trivial.
Como se tratava de uma questão cuja resposta exigia um conhecimento mínimo sobre o assunto, os estudantes que responderam, mesmo que parcialmente, o fizeram de acordo com o padrão de resposta.
As respostas dos estudantes foram, portanto, condizentes com padrão.
De maneira geral, as respostas foram razoáveis, não completas, mas adequadas.
Os estudantes mostraram que conhecem o assunto, porém não o dominam integralmente ou não se recordavam de todos os itens. 3.3.2.3 Análise da Questão Discursiva 4 do Componente de Conhecimento Específico A Tabela 3.13 contém as informações relativas à questão 4 do conjunto de questões discursivas do Componente de Conhecimento Específico.
O desempenho dos estudantes de todo o Brasil nesta questão foi inferior ao desempenho na questão de número 3, cuja nota média foi 38,5.
A média geral do Brasil na questão 4 foi 18,8, sendo a menor média registrada na região Nordeste ( 14,8 ), e a maior, na região Sul ( 20,9 ).
A nota máxima 95,0 foi atingida por pelo menos um aluno da região Sudeste.
A menor nota máxima foi obtida na região Norte ( 55,0 ).
A mediana do Brasil ( 10,0 ) foi a mesma nas regiões Norte e Sudeste.
A menor mediana foi 5,0, obtida na região Nordeste, e a maior foi 15,0, obtida nas regiões Sul e Centro-Oeste.
A nota mínima ( 0,0 ) foi obtida em todas as regiões do Brasil, sem exceção.
O Gráfico 3.20, representa a distribuição de notas na questão discursiva 4, no Componente de Conhecimento Específico.
A categoria “ em branco ” é a de maior frequência, sendo a moda.
Não se considerando esta categoria, o intervalo ( 10 ; 20 ] é o de maior frequência.
Na questão anterior, o intervalo de maior frequência, descartando -se as questões “ em branco ”, foi ( 40 ; 50 ], o que demonstra o menor desempenho na questão 4. 3.3.2.4 Comentários sobre a correção das respostas à Questão Discursiva 4 Trata -se de uma questão sobre conhecimento específico na área de sistemas de automação industrial, envolvendo itens relacionados a motores elétricos e seus acionamentos.
Simples, bem elaborada, exigindo que o estudante conheça as características dos motores normalmente empregados, as diferenças entre estes tipos de motores, as vantagens e desvantagens de se utilizar um determinado tipo, e os seus componentes de acionamento e controle.
A pontuação dos itens foi adequada à dificuldade.
A questão, como um todo, foi considerada de dificuldade média.
O padrão de resposta de todos os itens está de acordo com as referências bibliográficas clássicas e manuais técnicos da área.
Como se tratava de uma questão cuja resposta exigia conhecimento não superficial sobre o assunto, os estudantes que responderam, embora tenham sido poucos os que o fizeram, mesmo que parcialmente, responderam em coerência com o padrão de resposta.
As respostas daqueles que fizeram elaborações, de maneira geral, foram medianas, não completas, mas adequadas.
Alguns dos estudantes que responderam mostraram que conhecem o assunto, porém não o dominam integralmente ou não se recordavam de todos os itens.
No entanto, houve quantidade mais expressiva de notas baixas, entre zero e 30,0 pontos.
Há que se destacar o \textbf{elevado} número de respostas em branco. 3.3.2.5 Análise da Questão Discursiva 5 do Componente de Conhecimento Específico A Tabela 3.14 contém as informações relativas à questão 5 do conjunto de questões discursivas do Componente de Conhecimento Específico.
O desempenho dos estudantes nessa questão também foi inferior ao da questão 3.
Em relação à questão 4 as médias foram praticamente iguais, 18,8 na questão 4 e 18,2 na questão 5.
A maior média nesta questão foi registrada na região Sul ( 18,7 ), enquanto a menor média foi registrada na região Centro-Oeste ( 15,8 ).
Quanto à variabilidade das notas, o desvio padrão dos estudantes do Brasil, como um todo, foi 19,8.
Enquanto o maior desvio foi encontrado na região Sul ( 20,4 ), o menor foi encontrado na região Norte ( 18,1 ).
A nota máxima ( 100,0 ) foi alcançada na região Sul, enquanto a menor nota máxima foi registrada na região Norte ( 50,0 ).
A nota mínima foi zero para todas as regiões do Brasil.
A mediana foi 15,0 nas regiões Nordeste, Sudeste e Sul, a mesma do Brasil como um todo.
Na região Norte a mediana foi 12,5 e na região Centro-Oeste, 10,0.
O Gráfico 3.21 representa a distribuição das notas da questão discursiva 5 do Componente de Conhecimento Específico.
Destaca -se a grande frequência da categoria “ em branco ”, a maior dentre as três questões discursivas do Componente de Conhecimento Específico.
Como na questão 4, um máximo local ocorre no intervalo ( 10 ; 20 ].
Os coeficientes de assimetria das distribuições são positivos, tanto para o Brasil como um todo ( 0,94 ), representada no Gráfico 3.21, quanto de cada uma das Grandes Regiões ( variando entre 0,57 e 1,09 ). 3.3.2.6 Comentários sobre a correção das respostas à Questão Discursiva 5 Questão sobre conhecimento básico na área de sistemas de controle, tratando de itens relacionados a uma malha típica, aplicados a um sistema de aquecimento.
Simples, bem elaborada, exige que o estudante identifique e classifique as principais variáveis presentes na malha, associe os componentes do sistema com os elementos normalmente presentes em uma malha de controle, e especifique o transdutor a ser empregado em uma determinada situação.
A questão foi considerada com dificuldade média para difícil, de acordo com a avaliação da Banca de docentes corretores.
Como se tratava de uma questão cuja resposta exigia um conhecimento considerado fundamental sobre sua área de formação, alguns dos estudantes que responderam, mesmo que parcialmente, o fizeram de acordo com o padrão de resposta.
No entanto, o maior quantitativo de estudantes obteve notas baixas, com uma faixa expressiva variando entre zero e 30,0.
Em torno de 10 \% dos estudantes, por exemplo, tiveram notas entre zero e 10,0 pontos, e quase 15 \% tiveram notas na faixa entre 10,0 e 20,0 pontos.
Há que se destacar que esta foi a questão com maior quantidade de respostas deixadas em branco pelos estudantes, dentre as três questões discursivas de Conhecimento Específico. 3.3.3 Considerações Finais De acordo com a avaliação da Banca de docentes corretores, os alunos apresentaram dificuldade para resolver as questões discursivas do Componente de Conhecimento Específico, questões 3, 4 e 5.
As questões abordavam conteúdos importantes para o curso e com o resultado da prova não é possível afirmar que houve assimilação de tais conteúdos pelos estudantes.
Além disso, em função da grande quantidade de respostas deixadas em branco, as questões não contribuíram para avaliar estudantes e cursos.

% matched lemmas: creche, elevado, grandeza, isenção, vicioso
\end{textsample}
